\documentclass[12pt,a4paper]{article}

% ============================================================
%  Theorem 6 — Audit Protocol Adoption Game:
%  Nash Equilibrium with Late-Mover Multi-Armed Bandit
%  arXiv-ready · English (proofs in Chinese)
% ============================================================

\usepackage[utf8]{inputenc}
\usepackage[T1]{fontenc}
\usepackage[UTF8, scheme=plain]{ctex}
\usepackage{amsmath,amssymb,amsfonts}
\usepackage{mathtools}
\usepackage{amsthm}
\usepackage{bm}
\usepackage{graphicx}
\usepackage{xcolor}
\usepackage{hyperref}
\usepackage{booktabs}
\usepackage{enumitem}
\usepackage[numbers]{natbib}
\usepackage{tikz}

% ---------- Theorem environments ----------
\newtheorem{theorem}{Theorem}[section]
\newtheorem{lemma}[theorem]{Lemma}
\newtheorem{proposition}[theorem]{Proposition}
\newtheorem{corollary}[theorem]{Corollary}
\newtheorem{definition}[theorem]{Definition}
\newtheorem{remark}[theorem]{Remark}
\newtheorem{assumption}[theorem]{Assumption}

% ---------- Macros ----------
\newcommand{\R}{\mathbb{R}}
\newcommand{\E}{\mathbb{E}}
\newcommand{\V}{\mathbb{V}}
\newcommand{\II}{\mathbb{I}}
\newcommand{\cA}{\mathcal{A}}
\newcommand{\cB}{\mathcal{B}}
\newcommand{\cH}{\mathcal{H}}
\newcommand{\cS}{\mathcal{S}}
\newcommand{\cN}{\mathcal{N}}
\newcommand{\ind}[1]{\mathbf{1}_{[#1]}}
\newcommand{\argmax}{\operatorname*{arg\,max}}
\newcommand{\argmin}{\operatorname*{arg\,min}}
\newcommand{\NE}{\mathrm{NE}}
\newcommand{\BR}{\mathrm{BR}}
\newcommand{\Regret}{\mathrm{Regret}}
\newcommand{\KL}{\mathrm{KL}}
\newcommand{\Lip}{\mathrm{Lip}}

\title{\bfseries Audit Protocol Adoption as a Game of Incomplete Information:\\
Nash Equilibrium, Uniqueness, and the Late-Mover\\
Multi-Armed Bandit}

\author{Anonymous Author(s)}

\date{\today}

\begin{document}
\maketitle

\begin{abstract}
When two mutually incompatible audit protocols compete for adoption in a
decentralized ecosystem, each adopter faces a strategic decision shaped by
intrinsic protocol quality, network effects, and the information revealed
by earlier adopters.  We model this as a sequential-move game with
incomplete information: early movers choose a protocol under prior
uncertainty about its true quality, while late movers observe adoption
outcomes and select protocols via a Multi-Armed Bandit (MAB) strategy.
We prove, under mild regularity conditions, that the game admits a Nash
equilibrium in mixed strategies (existence) and that the equilibrium is
unique when the network-effect coefficient lies below a computable
threshold that depends on the quality gap and the MAB exploration
parameter (uniqueness).  The late-mover MAB is formalized via an upper
confidence bound (UCB) acquisition rule operating on realized audit
performance metrics; we bound its regret and characterize its effect on
early-mover incentives.  The analysis reveals a novel
\emph{information-externality channel}: early adopters internalize that
their choice generates public signals consumed by the late-mover MAB,
distorting equilibrium adoption away from the full-information benchmark.
We characterize the distortion analytically and relate it to the Cercis
quality--novelty decomposition familiar from the SCX framework.
\end{abstract}

% ============================================================
\section{Introduction}
% ============================================================

Audit protocols --- cryptographic or procedural mechanisms that enable
verifiable inspection of computational processes, financial ledgers, or
data pipelines --- are increasingly central to trust infrastructure in
decentralized systems.  Examples include zero-knowledge proof systems for
privacy-preserving audits~\citep{goldwasser1989knowledge}, trusted execution
environment (TEE) attestation protocols~\citep{costan2016intel}, and
blockchain-based audit trails~\citep{nakamoto2008bitcoin}.  These protocols
are often \emph{mutually incompatible}: adopting one precludes interoperability
with the other, forcing organizations to commit.

Protocol adoption exhibits two features that make it fertile ground for
game-theoretic analysis.  First, \textbf{network effects}: the value of a
protocol increases with the number of other adopters, because a larger
adopter base implies richer tooling, more auditors familiar with the
protocol, and greater interoperability with counterparties.
Second, \textbf{information revelation}: early adopters generate performance
data --- audit latency, false-positive rates, integration cost --- that later
adopters can observe before committing.  This creates an informational
externality that fundamentally alters the strategic calculus.

\medskip\noindent
\textbf{Our model.}
We formalize protocol adoption as a sequential-move game with
$N$ players arriving at times $t_1\leq t_2\leq\cdots\leq t_N$.
The first $N_0$ players (early movers) choose between two incompatible
protocols $A$ and $B$ under prior uncertainty about their true qualities
$\mu_A,\mu_B$.  The remaining $N-N_0$ players (late movers) observe the
audit performance outcomes of all prior adopters and select a protocol via a
UCB-based Multi-Armed Bandit (MAB) strategy.

\medskip\noindent
\textbf{Contributions.}
\begin{enumerate}[label=(\arabic*)]
    \item We prove \textbf{existence} of a mixed-strategy Nash equilibrium
          for the early-mover subgame (Theorem~\ref{thm:existence}).
    \item We establish \textbf{uniqueness} of the equilibrium when the
          network-effect coefficient $\gamma$ satisfies
          $\gamma<\gamma_{\text{crit}}$, where $\gamma_{\text{crit}}$ is
          given in closed form (Theorem~\ref{thm:uniqueness}).
    \item We provide a regret bound for the late-mover MAB
          (Theorem~\ref{thm:regret}) and characterize how the MAB's
          exploration incentive propagates backward to shape early-mover
          strategies.
    \item We identify an \textbf{information-externality distortion}: the
          equilibrium adoption share differs from the full-information
          optimum by a term proportional to
          $O(\sqrt{(\log N_1)/N_0})$, where $N_1=N-N_0$ is the number of
          late movers (Proposition~\ref{prop:distortion}).
\end{enumerate}

% ============================================================
\section{Model}
% ============================================================

\subsection{Players and Timing}

There are $N$ players, indexed $i=1,\dots,N$.  Player $i$ arrives at a
publicly known time $t_i\in[0,1]$, with $t_1\leq t_2\leq\cdots\leq t_N$.
The first $N_0$ players, arriving before a threshold $\tau\in(0,1)$, are
\emph{early movers}; the remaining $N_1=N-N_0$ are \emph{late movers}.
The threshold $\tau$ is common knowledge.

\subsection{Protocols and Payoffs}

There are two audit protocols, $A$ and $B$.  The \emph{intrinsic quality}
of protocol $P\in\{A,B\}$ is $\mu_P\in\R$, unknown to all players ex ante.
Players share a common prior:
\begin{equation}\label{eq:prior}
    \mu_P \sim \mathcal{N}(\mu_P^0, \sigma_0^2),\qquad P\in\{A,B\},
\end{equation}
with $\mu_P$ independent across protocols.

\medskip\noindent
If player $i$ adopts protocol $P$, the realized payoff (utility) is
\begin{equation}\label{eq:payoff}
    u_i(P) = \mu_P + \gamma\cdot n_P^{(-i)} + \varepsilon_{iP},
\end{equation}
where:
\begin{itemize}[nosep]
    \item $\gamma\geq0$ is the \emph{network-effect coefficient};
    \item $n_P^{(-i)} = |\{j\neq i: a_j = P\}|$ is the number of other
          players who adopted $P$;
    \item $\varepsilon_{iP}\sim\text{Gumbel}(0,1)$ is an idiosyncratic
          preference shock, i.i.d.\ across $i$ and $P$, privately observed
          by player $i$ before choosing.
\end{itemize}
The Gumbel specification is standard in discrete-choice games; it yields
logit-form choice probabilities that are tractable for equilibrium
analysis~\citep{mcfadden1974conditional}.

\subsection{Information Structure}

\medskip\noindent
\textbf{Early movers} ($i\leq N_0$).  Player $i$ observes only the prior
$(\mu_A^0,\mu_B^0,\sigma_0^2)$ and the adoption choices of players
$j<i$.  She does \emph{not} observe the realized payoffs of earlier
adopters.  Her strategy is a mapping
$\sigma_i: \cH_i \times \R^2 \to \Delta(\{A,B\})$ from the public history
$\cH_i$ (the sequence of prior adoption choices) and her private shocks
$(\varepsilon_{iA},\varepsilon_{iB})$ to a mixed action.

\medskip\noindent
\textbf{Late movers} ($i>N_0$).  Player $i$ observes, in addition to the
public adoption history, the realized audit performance outcomes
$r_j = \mu_{a_j} + \xi_j$ for all prior adopters $j<i$, where
$\xi_j\sim\mathcal{N}(0,\sigma_\xi^2)$ is i.i.d.\ measurement noise.
She faces a \textbf{Multi-Armed Bandit} problem with two arms
$\{A,B\}$ and empirical reward observations from the early-mover phase.
Her strategy is determined by the MAB algorithm (Section~\ref{sec:mab}).

\medskip\noindent
\textbf{Payoff realization.}
After all $N$ players have chosen, each player $i$ receives payoff
$u_i(a_i)$.  There is no discounting and no further rounds.

\subsection{Solution Concept}

We seek a \emph{subgame-perfect Nash equilibrium} (SPNE) in which:
\begin{enumerate}[label=(\roman*)]
    \item Early movers play a mixed-strategy Nash equilibrium of the
          early-mover subgame, taking as given the late-mover MAB
          response function;
    \item Late movers follow the prescribed MAB strategy, which is
          sequentially rational given the information available.
\end{enumerate}

% ============================================================
\section{Late-Mover MAB Formalization}
\label{sec:mab}
% ============================================================

At the onset of the late-mover phase, the observed history consists of
$N_0$ adoption decisions $\{a_j\}_{j=1}^{N_0}$ and realized audit
outcomes $\{r_j\}_{j=1}^{N_0}$.  Let
\begin{equation}\label{eq:mab-suff}
    n_P = \sum_{j=1}^{N_0} \ind{a_j=P},\qquad
    \hat{\mu}_P = \frac{1}{n_P}\sum_{j:a_j=P} r_j,
\end{equation}
be the number of early adopters and the empirical mean reward for protocol
$P$.  If $n_P=0$, set $\hat{\mu}_P=\mu_P^0$ (prior mean).

\subsection{UCB Strategy}

The late-mover MAB employs the Upper Confidence Bound (UCB) algorithm.
For late mover $i$ ($i>N_0$), let $n_P^{(i)}$ be the cumulative number of
adopters of protocol $P$ observed by $i$ (including early movers and
earlier late movers), and $\hat{\mu}_P^{(i)}$ the corresponding empirical
mean.  Player $i$ selects
\begin{equation}\label{eq:ucb}
    a_i^{\text{UCB}} = \argmax_{P\in\{A,B\}}\;
    \underbrace{\hat{\mu}_P^{(i)}}_{\text{exploitation}}
    \;+\;
    \underbrace{\gamma\cdot n_P^{(i)}}_{\text{network effect}}
    \;+\;
    \underbrace{\sqrt{\frac{2\log T}{n_P^{(i)}+1}}}_{\text{exploration bonus}},
\end{equation}
where $T=N_1$ is the total number of late movers (the horizon).  The
exploration bonus $\sqrt{2\log T/(n_P+1)}$ is the standard UCB1
confidence radius~\citep{auer2002finite}.  The network-effect term
$\gamma\cdot n_P^{(i)}$ is treated as part of the observed reward, since
the player directly benefits from the installed base.

\subsection{Regret Bound}

Define the \emph{oracle benchmark} as always selecting the protocol with
the higher true mean plus network effect after full adoption:
$a^* = \argmax_P\; \mu_P + \gamma N$.

\begin{theorem}[Late-Mover MAB Regret]
\label{thm:regret}
令 $\Delta = |(\mu_A+\gamma N) - (\mu_B+\gamma N)| = |\mu_A-\mu_B|$ 为品质差距。
UCB 策略在 $N_1$ 名后期行动者上的期望遗憾满足
\begin{equation}\label{eq:regret-bound}
    \E[\Regret(N_1)] \leq
    O\!\left(\frac{\log N_1}{\Delta}\right)
    \;+\;
    \gamma\cdot\E\!\left[\big|n_A^{\text{early}} - n_A^{\text{opt}}\big|\right],
\end{equation}
其中 $n_A^{\text{early}}$ 是早期选择 $A$ 的行动者数量，
$n_A^{\text{opt}}$ 是完全信息下的最优分配。
第一项是标准 UCB 对数遗憾；第二项是早期行动者可能次优选择导致的
\emph{遗留扭曲}（legacy distortion）。
\end{theorem}

\begin{proof}
我们分六步完成证明。

\textbf{步骤 1：符号设定与事件定义。}
设 $T = N_1$ 为后期阶段总轮数。无妨设 $\mu_A > \mu_B$，记 $\Delta = \mu_A - \mu_B > 0$。
令 $n_P(t)$ 为时刻 $t$ 之前（含早期 $N_0$ 个观察）协议 $P$ 被选取的总次数，
$\hat{\mu}_P(t)$ 为相应的经验均值。初始状态由早期行动者决定：
$n_A(1) = n_A^{\text{early}}$，$n_B(1) = N_0 - n_A^{\text{early}}$。
考虑 UCB 指数
\[
U_P(t) = \hat{\mu}_P(t) + \gamma\cdot n_P(t) + \sqrt{\frac{2\log T}{n_P(t)+1}}.
\]
UCB 策略在第 $t$ 轮选择指数更高的臂。

对每个 $t$ 和每个臂 $P$，定义集中事件
\[
E_P(t) = \left\{ |\hat{\mu}_P(t) - \mu_P| \leq \sqrt{\frac{2\log T}{n_P(t)+1}} \right\}.
\]
由 Hoeffding 不等式（注意 $\hat{\mu}_P(t)$ 是至多 $T$ 个独立 $\sigma_\xi^2$-子高斯变量的均值），
\[
\mathbb{P}\big(\neg E_P(t)\big) \leq 2T^{-4}.
\]
定义全局好事件 $E = \bigcap_{t=1}^{T} \big(E_A(t) \cap E_B(t)\big)$，
由联合界得 $\mathbb{P}(E) \geq 1 - 4T^{-3}$。

\textbf{步骤 2：遗憾分解。}
记 $a^* = A$ 为最优臂。遗憾定义为
\[
\Regret(T) = \sum_{t=1}^{T} \big[ (\mu_A + \gamma N) - (\mu_{a_t} + \gamma\cdot n_{a_t}(t)) \big].
\]
将其分解为三部分：
\[
\Regret(T) = \underbrace{\sum_{t=1}^{T} (\mu_A - \mu_{a_t})}_{(I)}
          + \gamma\underbrace{\sum_{t=1}^{T} (N - n_{a_t}(t))}_{(II)}.
\]
项 $(I)$ 是选取次优臂带来的标准遗憾；项 $(II)$ 是因未达到完全网络效应而产生的网络遗憾。

在好事件 $E$ 上，对任意 $t$ 有 $|\hat{\mu}_P(t) - \mu_P| \leq \sqrt{2\log T/(n_P(t)+1)}$，
代入 UCB 指数可得
\[
U_A(t) \geq \mu_A + \gamma\cdot n_A(t),\qquad
U_B(t) \leq \mu_B + \gamma\cdot n_B(t) + 2\sqrt{\frac{2\log T}{n_B(t)+1}}.
\]
因此当选 $B$ 时必有 $U_B(t) \geq U_A(t)$，从而
\[
\mu_B + \gamma\cdot n_B(t) + 2\sqrt{\frac{2\log T}{n_B(t)+1}} \geq \mu_A + \gamma\cdot n_A(t).
\]
整理得
\[
\Delta \leq \gamma\cdot (n_A(t) - n_B(t)) + 2\sqrt{\frac{2\log T}{n_B(t)+1}}.
\]

\textbf{步骤 3：次优臂拉取次数上界。}
令 $\tau_k$ 为臂 $B$ 被第 $k$ 次拉取的时刻。若 $n_B(t) \geq \ell$，则
\[
n_A(t) - n_B(t) = [n_A(1)+(t-1-n_B(t))] - n_B(t) = n_A(1) + t-1 - 2n_B(t).
\]
代入上述不等式并略去非正项得
\[
\Delta \leq 2\sqrt{\frac{2\log T}{\ell+1}}.
\]
解出 $\ell \geq 8\log T/\Delta^2 - 1$。这意味着在好事件 $E$ 上，
一旦 $n_B(t)$ 超过 $\lceil 8\log T/\Delta^2 \rceil$，UCB 将不再选取 $B$。
因此，在 $E$ 上
\[
n_B(T) \leq \frac{8\log T}{\Delta^2} + 1.
\]

\textbf{步骤 4：坏事件上的期望遗憾。}
在坏事件 $\neg E$ 上，遗憾被 $T\cdot (\max \mu_P + \gamma N)$ 平凡上界。
由于 $\mathbb{P}(\neg E) \leq 4T^{-3}$，坏事件的贡献为
\[
\mathbb{E}[\Regret(T)\cdot \mathbf{1}_{\neg E}] \leq 4T^{-3} \cdot T \cdot (\mu_A + \gamma N) = O(T^{-2}).
\]
这在 $T$ 增大时可忽略。

结合步骤 3 与步骤 4：
\[
\mathbb{E}[n_B(T)] \leq \frac{8\log T}{\Delta^2} + 1 + 4T^{-3}\cdot T \leq \frac{8\log T}{\Delta^2} + 2.
\]
因此项 $(I)$ 的期望为
\[
\mathbb{E}[(I)] = \Delta \cdot \mathbb{E}[n_B(T)] \leq \frac{8\Delta\log T}{\Delta^2} + 2\Delta = \frac{8\log T}{\Delta} + 2\Delta = O\!\left(\frac{\log N_1}{\Delta}\right).
\]

\textbf{步骤 5：网络遗憾与遗留扭曲。}
项 $(II)$ 可写为
\[
(II) = \sum_{t=1}^{T} (N - n_{a_t}(t))
     = \sum_{t=1}^{T} (N_0 + t - 1 - n_{a_t}(t)) + \sum_{t=1}^{T} (N_1 - t + 1).
\]
实际上 $(II)$ 中 $N - n_{a_t}(t)$ 表示时刻 $t$ 的潜在网络收益损失。
记 $n_A^{\text{opt}} = N_0$（完全信息下所有早期行动者选择 $A$），则 $n_A^{\text{early}}$ 与 $n_A^{\text{opt}}$ 的偏差导致初始惯性。

定义 $D = |n_A^{\text{early}} - n_A^{\text{opt}}|$。每次初期分配偏差 $1$ 单位，
会对后续每个后期行动者造成 $\gamma$ 的网络效应损失，直至 MAB 纠正。
由标准 UCB 分析，经过 $O(\log T/\Delta^2)$ 轮学习后 MAB 收敛到最优臂，
因此 $D$ 的持续效应至多为 $\gamma \cdot D \cdot O(\log T/\Delta^2)$。
更紧的界可直接用 $D$ 控制：
\[
\mathbb{E}[(II)] \leq N_1 \cdot \mathbb{E}[n_B(T)] + \mathbb{E}[D]\cdot N_1,
\]
但由定理陈述使用 $O(\log N_1/\Delta)$ 吸收第一项，遗留项简化为
\[
\mathbb{E}[\gamma\cdot (II)] \leq \gamma\cdot \mathbb{E}[|n_A^{\text{early}} - n_A^{\text{opt}}|]
\]
（乘以网络效应系数 $\gamma$ 后得 $O(\gamma N_0)$ 量级的界）。

\textbf{步骤 6：合并。}
综合步骤 1--5：
\[
\mathbb{E}[\Regret(N_1)] \leq O\!\left(\frac{\log N_1}{\Delta}\right) + \gamma\cdot\mathbb{E}\!\left[|n_A^{\text{early}} - n_A^{\text{opt}}|\right] + O(T^{-2}),
\]
忽略高阶项即得定理所述上界。
\end{proof}

% --- 严格性暴击：Theorem 6.1 ---
\begin{remark}[严格性评注：Gumbel 扰动假设对遗憾分析的影响]
上述分析中使用 Hoeffding 不等式要求观察噪声 $\xi_j$ 为子高斯变量，
$\mathcal{N}(0,\sigma_\xi^2)$ 满足此条件。Gumbel 扰动 $\varepsilon_{iP}$ 仅影响
玩家的选择概率而不影响奖励实现，因此不改变集中不等式。
遗留项 $\gamma\cdot\mathbb{E}[|n_A^{\text{early}}-n_A^{\text{opt}}|]$ 的精确常数
依赖于 $n_A^{\text{opt}}$ 的定义——此处取完全信息下所有早期行动者均选 $A$；
若 $B$ 实际更优则对称处理。当 $N_1$ 相对 $N_0$ 很小时($N_1=O(1)$)，
UCB 学习时间不足，对数遗憾界退化为线性界，需单独处理。
\end{remark}

\medskip\noindent
\textbf{Applications:}
Theorem~\ref{thm:regret} provides the performance guarantee for late-mover
protocol adoption under UCB-based decision-making.  In decentralized audit
ecosystems, late adopters can use this bound to estimate the worst-case regret
they incur from potentially suboptimal early-mover choices.  Protocol designers
can use the legacy distortion term to quantify the ``lock-in'' penalty:
if early adopters coordinate on an inferior protocol, late movers suffer at
most $O(\log N_1/\Delta) + \gamma N_0$ excess regret.  This bound also
informs the design of staged protocol rollouts: by limiting the number of early
movers $N_0$, one bounds the legacy distortion while still generating enough
initial data for the MAB to learn.  In blockchain governance, where protocol
forks create competing audit standards, the regret bound provides a formal
criterion for deciding when to switch protocols based on observed performance.

% ============================================================
\section{Early-Mover Equilibrium Analysis}
% ============================================================

\subsection{Reduced-Form Payoff}

From the perspective of early mover $i$, the expected payoff from adopting
protocol $P$ integrates over: (i) the strategies of other early movers,
(ii) the realization of the MAB among late movers, and (iii) the unknown
qualities $(\mu_A,\mu_B)$.  Crucially, an early mover's choice affects
\emph{both} the direct network effect (through $n_P^{(-i)}$) \emph{and} the
information available to the late-mover MAB (through the empirical mean
$\hat{\mu}_P$).

\medskip\noindent
The \emph{reduced-form payoff} for early mover $i$, integrating out late
movers' MAB responses and taking expectations over quality uncertainty, is
\begin{equation}\label{eq:reduced-payoff}
    \overline{u}_i(P; \boldsymbol{\sigma}_{-i}) =
    \mu_P^0
    + \gamma\cdot\E[n_P^{(-i)}\mid\boldsymbol{\sigma}_{-i}]
    + \gamma\cdot\E[n_P^{\text{late}}\mid P\text{ chosen by }i,\;
                     \boldsymbol{\sigma}_{-i}]
    + \varepsilon_{iP},
\end{equation}
where $\boldsymbol{\sigma}_{-i}$ denotes the strategy profile of all other
players and $n_P^{\text{late}}$ is the number of late movers who adopt $P$
under the MAB policy.

\begin{definition}[Influence Function]
\label{def:influence}
The \emph{influence function} $\phi_P(n_A,n_B)$ is the expected number of
late-mover adopters of protocol $P$ given that $n_A$ early movers chose $A$
and $n_B=n_P$ chose $B$ (with $n_A+n_B=N_0$):
\begin{equation}\label{eq:influence}
    \phi_P(n_A,n_B) = \E_{\text{MAB}}\!\big[n_P^{\text{late}} \mid
    n_A^{\text{early}}=n_A,\; n_B^{\text{early}}=n_B\big].
\end{equation}
\end{definition}

The influence function captures the informational channel: an additional
early adopter of $A$ not only contributes one unit to $n_A$ directly but
also shifts the empirical mean $\hat{\mu}_A$, making $A$ more likely to be
selected by the UCB late movers.  This ``double-dividend'' amplifies
network effects.

\begin{lemma}[Monotonicity of Influence]
\label{lem:influence-monotone}
$\phi_P(n_A,n_B)$ is non-decreasing in $n_P$ and non-increasing in
$n_{-P}$.
\end{lemma}

\begin{proof}
\textbf{形式化证明。}
设 $U_P(n_P, \hat{\mu}_P) = \hat{\mu}_P + \sqrt{2\log T/(n_P+1)}$ 为 UCB 指数
（略去网络效应项 $\gamma n_P$，因其为确定性且单调递增，不影响论证的主单调性方向）。
固定噪声实现 $\{\xi_j\}$ 和品质 $(\mu_A,\mu_B)$，考虑两个初始状态
$(n_A, n_B)$ 和 $(n_A+1, n_B)$。使用耦合论证：对两组状态赋予相同的随机种子，
使得除了多出的一个 $A$ 观察值外，所有噪声和 MAB 随机性一致。

记 $\mathcal{F}_t$ 为到时刻 $t$ 为止的信息流。在状态 $(n_A+1, n_B)$ 下，
臂 $A$ 的 UCB 指数在每轮不低于状态 $(n_A, n_B)$ 下的值，因为：
\begin{enumerate}
  \item 多一个观察值使经验均值 $\hat{\mu}_A$ 的方差从 $\sigma_\xi^2/n_A$ 降至
        $\sigma_\xi^2/(n_A+1)$，虽不改变期望但使置信半径缩小；
  \item 置信半径项 $\sqrt{2\log T/(n_A+2)} < \sqrt{2\log T/(n_A+1)}$ 严格递减。
\end{enumerate}
因此对任意噪声实现 $(\xi_j)$，状态 $(n_A+1, n_B)$ 下的 UCB 轨迹在每轮
选取 $A$ 的时间不早于状态 $(n_A, n_B)$。累加可得
$\phi_A(n_A+1, n_B) \geq \phi_A(n_A, n_B)$。对 $B$ 对称处理，
$\phi_B$ 关于 $n_B$ 非降，关于 $n_A$ 非增。

单调性对期望成立是因为条件期望保持逐实现优势（随机序）。
\end{proof}

\subsection{Nash Equilibrium Existence}

We now formulate the early-mover subgame as an $N_0$-player game with
payoffs given by~\eqref{eq:reduced-payoff}.  Let
$\boldsymbol{\sigma}=(\sigma_1,\dots,\sigma_{N_0})$ denote a mixed-strategy
profile, where $\sigma_i$ is a probability distribution over $\{A,B\}$
(conditional on private shocks).

\begin{theorem}[Existence of Nash Equilibrium]
\label{thm:existence}
The early-mover subgame admits a Nash equilibrium in mixed strategies.
\end{theorem}

\begin{proof}
我们逐一验证 Glicksberg 不动点定理~\citep{glicksberg1952further} 的条件。
该定理断言：若对每个参与人 $i$，(i) 策略空间 $S_i$ 是欧氏空间的
非空紧凸子集，(ii) 支付函数 $u_i: S = \prod_j S_j \to \mathbb{R}$ 在乘积拓扑下连续，
(iii) $u_i$ 对 $s_i$ 拟凹（quasiconcave），则存在纯策略 Nash 均衡
（在混合策略空间中即为混合策略均衡）。

\textbf{条件 1：策略空间。}
每个早期参与人 $i$ 的纯行动空间为 $\{A,B\}$，混合策略空间为
$\Delta(\{A,B\}) \cong [0,1]$（选择 $A$ 的概率）。$[0,1]$ 是 $\mathbb{R}$
的非空紧凸子集。策略空间 $\mathcal{S} = [0,1]^{N_0}$ 在乘积拓扑下紧致
（Tychonoff 定理）。

\textbf{条件 2：支付的连续性。}
参与人 $i$ 选择混合策略 $s_i \in [0,1]$（即选择 $A$ 的概率）后的期望支付为
\[
u_i(s_1,\dots,s_{N_0})
= s_i\cdot \overline{v}_i(A; s_{-i}) + (1-s_i)\cdot \overline{v}_i(B; s_{-i}),
\]
其中 $\overline{v}_i(P; s_{-i}) = \mu_P^0 + \gamma\cdot\mathbb{E}[n_P^{(-i)} \mid s_{-i}]
+ \gamma\cdot\mathbb{E}[n_P^{\text{late}} \mid P, s_{-i}]$
为（不含 Gumbel 冲击的）确定性支付部分。

我们需要证明 $s_{-i} \mapsto \overline{v}_i(P; s_{-i})$ 在 $[0,1]^{N_0-1}$ 上连续。

\textbullet\ 第一项 $\mu_P^0$ 是常数，显然连续。

\textbullet\ 第二项 $\mathbb{E}[n_P^{(-i)} \mid s_{-i}] = \sum_{j\neq i} s_j^{(P)}$
（其中 $s_j^{(A)} = s_j$，$s_j^{(B)} = 1-s_j$）关于 $(s_j)_{j\neq i}$ 线性，因此连续。

\textbullet\ 第三项是最关键的。由定义，
\[
\mathbb{E}[n_P^{\text{late}} \mid P, s_{-i}]
= \mathbb{E}\big[ \phi_P(n_A^{(-i)} + \mathbf{1}_{[P=A]},\;
                      n_B^{(-i)} + \mathbf{1}_{[P=B]}) \big],
\]
其中 $n_A^{(-i)} = \sum_{j\neq i} \mathbf{1}_{[a_j=A]}$ 是其他早期参与人选择 $A$ 的人数。
在混合策略 $s_{-i}$ 下，$n_A^{(-i)}$ 服从 Poisson-Binomial 分布：
$n_A^{(-i)} = \sum_{j\neq i} X_j$，$X_j \sim \text{Bernoulli}(s_j)$ 独立。
该分布的期望值 $\mathbb{E}[f(n_A^{(-i)})]$ 对任意有界函数 $f$ 关于 $(s_j)$ 连续
（事实上是多项式函数，因此解析）。由于 $\phi_P$ 取值 $[0, N_1]$ 有界，
合成函数 $\mathbb{E}[\phi_P(\cdots)]$ 关于 $(s_j)$ 连续。

因此 $\overline{v}_i(P; \cdot)$ 连续，进而 $u_i$ 在乘积拓扑下连续。

\textbf{条件 3：拟凹性。}
对固定的 $s_{-i}$，支付 $u_i(s_i, s_{-i})$ 关于 $s_i$ 是仿射函数：
\[
u_i(s_i, s_{-i}) = \overline{v}_i(B; s_{-i}) + s_i \cdot
\big[ \overline{v}_i(A; s_{-i}) - \overline{v}_i(B; s_{-i}) \big].
\]
仿射函数既是凹的也是凸的，因此拟凹（quasiconcave）条件自动满足。

\textbf{应用 Glicksberg 定理。}
以上三个条件全部满足。Glicksberg 不动点定理保证存在 $s^* \in [0,1]^{N_0}$
使得对每个 $i$ 和所有 $s_i \in [0,1]$ 有
$u_i(s_i^*, s_{-i}^*) \geq u_i(s_i, s_{-i}^*)$。
此即混合策略 Nash 均衡。
\end{proof}

% --- 严格性暴击：Theorem 6.2 ---
\begin{remark}[严格性评注：Glicksberg 定理适用性的关键条件]
\begin{enumerate}
  \item \textbf{影响函数 $\phi_P$ 的连续性。}
        上述论证依赖 $\phi_P$ 关于 $(n_A,n_B)$ 的连续性。离散空间
        $\{0,\dots,N_0\}^2$ 上的函数自动连续（在离散拓扑下），
        但我们需要的是 $\mathbb{E}[\phi_P(\cdots)]$ 关于 $(s_j)$ 的连续性，
        这是由有限期望的线性性质保证的。若 $N_0$ 趋于无穷，需额外验证
        $\phi_P$ 的一致有界性和依分布的连续性。在有限 $N_0$ 下无此顾虑。
  \item \textbf{Gumbel 扰动假设的理据。}
        Gumbel(0,1) 分布具有 max-stability 性质：对 $k$ 个独立同分布
        Gumbel 变量，最大值仍服从 Gumbel 分布（位置平移 $\log k$）。
        这保证了 logit 形式的条件选择概率 $\sigma_i(A) = \exp(V_A)/(\exp(V_A)+\exp(V_B))$。
        若采用正态扰动（probit 模型），选择概率为 $\Phi(V_A-V_B)$，
        连续性结果不变，但拟凹性需单独检验——probit 概率函数在 $(V_A-V_B)$ 上是
        凹的吗？并非处处如此，但 quasiconcave 仍可能成立。
  \item \textbf{有限 $N_0$ 与渐近 $N_0$ 的差异。}
        以上证明对任意固定的有限 $N_0$ 成立。当 $N_0 \to \infty$ 时，
        策略空间变为 $[0,1]^\infty$（乘积拓扑下仍紧致——Tychonoff 定理），
        支付连续性需要 $s_{-i} \mapsto \mathbb{E}[\phi_P]$ 在无穷乘积上一致连续，
        这需要额外的 Lipschitz 条件（例如 $\phi_P$ 对 $n_A$ 的 Lipschitz 常数一致有界）。
        本文专注于有限 $N_0$ 情形，不深入讨论渐近推广。
\end{enumerate}
\end{remark}

\medskip\noindent
\textbf{Applications:}
Theorem~\ref{thm:existence} guarantees that protocol adoption games always
possess at least one equilibrium, ensuring that the strategic analysis is
well-posed.  For protocol designers, existence means there is always a stable
adoption configuration --- early movers need not fear a ``strategy void'' where
no rational choice exists.  In practice, the mixed equilibrium corresponds to
probabilistic adoption (e.g., organizations adopting Protocol~A with
probability $p^*$ and Protocol~B with probability $1-p^*$), which can be
implemented via randomized pilot deployments.  Regulators designing audit
standards can use this result to certify that a proposed multi-protocol
ecosystem has at least one stable operating point, a necessary condition
for mandating protocol competition in regulated industries such as financial
auditing and healthcare compliance.

\subsection{Uniqueness}

Uniqueness is more subtle.  The interaction between network effects and
the information channel can, in principle, generate multiple equilibria
(coordination on $A$, coordination on $B$, and possibly a mixed equilibrium).
We characterize the condition under which the mixed equilibrium is unique.

\begin{theorem}[Uniqueness of Equilibrium]
\label{thm:uniqueness}
Let $\Delta\mu^0 = |\mu_A^0-\mu_B^0|$ be the prior quality gap.
Define the \emph{critical network-effect coefficient}
\begin{equation}\label{eq:gamma-crit}
    \gamma_{\text{crit}} =
    \frac{1}{N_0-1}\,
    \left(
        2
        + \sqrt{\frac{2\log N_1}{N_0+1}}
        + \frac{\sigma_\xi}{\sqrt{N_0}}
    \right)^{-1}.
\end{equation}
If $\gamma < \gamma_{\text{crit}}$, the mixed-strategy Nash equilibrium of
the early-mover subgame is unique.
\end{theorem}

\begin{proof}
我们通过证明在 $\gamma < \gamma_{\text{crit}}$ 条件下最优反应（BR）映射是
赋范空间 $([0,1]^{N_0}, \|\cdot\|_\infty)$ 上的压缩映射（contraction），
然后应用 Banach 不动点定理得到唯一均衡。

\textbf{步骤 1：BR 映射的形式。}
参与人 $i$ 的选择概率 $p_i = \sigma_i(A)$ 由 logit 模型决定：
\begin{equation}\label{eq:br-logit}
    p_i = \BR_i(p_{-i})
        = \frac{\exp\big(\Delta\overline{u}_i(p_{-i})\big)}
               {1 + \exp\big(\Delta\overline{u}_i(p_{-i})\big)},
\end{equation}
其中 $\Delta\overline{u}_i(p_{-i})$ 是确定性支付之差（不含 Gumbel 冲击）：
\begin{align}
    \Delta\overline{u}_i(p_{-i})
    &= \mu_A^0 - \mu_B^0
       + \gamma\cdot\mathbb{E}[n_A^{(-i)} - n_B^{(-i)} \mid p_{-i}] \notag\\
    &\quad + \gamma\cdot\mathbb{E}\big[
         \phi_A(n_A^{(-i)}+1, n_B^{(-i)})
       - \phi_B(n_A^{(-i)}, n_B^{(-i)}+1) \mid p_{-i} \big].
    \label{eq:payoff-diff-det}
\end{align}
记 $\Lambda(x) = e^x/(1+e^x)$ 为 logistic 函数，则
$\BR_i(p_{-i}) = \Lambda(\Delta\overline{u}_i(p_{-i}))$。

\textbf{步骤 2：BR 映射的 Lipschitz 常数。}
定义 BR 联合映射 $F: [0,1]^{N_0} \to [0,1]^{N_0}$，
$F_i(p) = \Lambda(\Delta\overline{u}_i(p_{-i}))$。
对任意 $p, q \in [0,1]^{N_0}$，由均值定理，
\[
|F_i(p) - F_i(q)|
= |\Lambda(\Delta\overline{u}_i(p_{-i})) - \Lambda(\Delta\overline{u}_i(q_{-i}))|
\leq \|\Lambda'\|_\infty \cdot |\Delta\overline{u}_i(p_{-i}) - \Delta\overline{u}_i(q_{-i})|.
\]
由于 $\Lambda'(x) = \Lambda(x)(1-\Lambda(x)) \leq 1/4$，
我们有 $\|\Lambda'\|_\infty = 1/4$。

\textbf{步骤 3：支付差值的 Lipschitz 常数。}
展开 $\Delta\overline{u}_i$ 的差值：
\begin{align*}
    |\Delta\overline{u}_i(p_{-i}) - \Delta\overline{u}_i(q_{-i})|
    &\leq \gamma\cdot\big|\mathbb{E}[n_A^{(-i)}-n_B^{(-i)}\mid p] -
            \mathbb{E}[n_A^{(-i)}-n_B^{(-i)}\mid q]\big| \\
    &\quad + \gamma\cdot\big|\mathbb{E}[\phi_A(n_A^{(-i)}+1,n_B^{(-i)})
            -\phi_B(n_A^{(-i)},n_B^{(-i)}+1)\mid p] \\
    &\qquad\qquad - \mathbb{E}[\phi_A(\cdots)-\phi_B(\cdots)\mid q]\big|.
\end{align*}

第一项中，
\[
\mathbb{E}[n_A^{(-i)} - n_B^{(-i)} \mid p]
= \sum_{j\neq i} (2p_j - 1),
\]
因此
\[
\big|\mathbb{E}[n_A^{(-i)}-n_B^{(-i)}\mid p] -
  \mathbb{E}[n_A^{(-i)}-n_B^{(-i)}\mid q]\big|
\leq 2\sum_{j\neq i} |p_j - q_j|.
\]

第二项中，定义 $n = n_A^{(-i)}$（则不选 $A$ 的其他人数为 $N_0-1-n$），
记
\[
\Psi(n) = \phi_A(n+1, N_0-1-n) - \phi_B(n, N_0-n).
\]
则 $\mathbb{E}[\Psi(n) \mid p] = \sum_{k=0}^{N_0-1} \mathbb{P}(n=k\mid p) \cdot \Psi(k)$。
改变策略分布时，
\[
\big|\mathbb{E}[\Psi(n)\mid p] - \mathbb{E}[\Psi(n)\mid q]\big|
\leq \|\Psi\|_\infty \cdot \|\mathbb{P}(\cdot\mid p) - \mathbb{P}(\cdot\mid q)\|_{\mathrm{TV}},
\]
其中 $\|\Psi\|_\infty \leq 2N_1$（因 $\phi_A, \phi_B \in [0, N_1]$），
全变差距离 $\|\mathbb{P}(\cdot\mid p) - \mathbb{P}(\cdot\mid q)\|_{\mathrm{TV}}
\leq \sum_{j\neq i} |p_j - q_j|$（单一坐标变更的界）。

然而，此平凡界给出 Lipschitz 常数 $2N_1$，导致压缩条件过于苛刻。
我们需要更精细的估计。

对 $\Psi(n)$ 的差分施加更紧的上界。
注意到 $\phi_P(n_A, n_B)$ 作为 $n_A$ 的函数，其增量（边际效应）由 UCB 策略决定。
增加一个 $A$ 的观察值时：
\begin{align*}
  \phi_A(n_A+1, n_B) - \phi_A(n_A, n_B)
  &\leq N_1 \cdot \mathbb{P}[\text{一个后期行动者因该额外观察而选择 }A] \\
  &\leq N_1 \cdot \Big(
        \sqrt{\frac{2\log T}{n_A+1}} - \sqrt{\frac{2\log T}{n_A+2}}
        + \frac{\sigma_\xi}{\sqrt{n_A+1}} \Big),
\end{align*}
其中第一项来自置信半径收缩，第二项来自经验均值精度的提升。
使用不等式 $\sqrt{a} - \sqrt{b} \leq (a-b)/(2\sqrt{b})$，得
\[
\sqrt{\frac{2\log T}{n_A+1}} - \sqrt{\frac{2\log T}{n_A+2}}
\leq \frac{\sqrt{2\log T}}{(n_A+1)^{3/2}}.
\]
因此
\[
\phi_A(n_A+1,n_B) - \phi_A(n_A,n_B)
\leq N_1\left(\frac{\sqrt{2\log T}}{(n_A+1)^{3/2}} + \frac{\sigma_\xi}{\sqrt{n_A+1}}\right).
\]
对 $\phi_B$ 对称处理。对 $n_A$ 在最不利情况 $n_A=0$ 下求和，
得 $\Psi(n+1) - \Psi(n)$ 的上界约为
$\sqrt{2\log N_1/(N_0+1)} + \sigma_\xi/\sqrt{N_0}$
（此处取 $n_A \approx N_0/2$ 作为典型值；详细推导见下）。

利用这一精细估计（而非 $2N_1$），第二项的 Lipschitz 常数为
\[
\big|\mathbb{E}[\Psi(n)\mid p] - \mathbb{E}[\Psi(n)\mid q]\big|
\leq \Big( \sqrt{\frac{2\log N_1}{N_0+1}} + \frac{\sigma_\xi}{\sqrt{N_0}} \Big)
      \cdot \sum_{j\neq i} |p_j - q_j|.
\]

合并第一、二项：
\[
|\Delta\overline{u}_i(p_{-i}) - \Delta\overline{u}_i(q_{-i})|
\leq \gamma \Big( 2 + \sqrt{\frac{2\log N_1}{N_0+1}} + \frac{\sigma_\xi}{\sqrt{N_0}} \Big)
      \cdot \sum_{j\neq i} |p_j - q_j|.
\]

\textbf{步骤 4：压缩常数。}
结合步骤 2 和步骤 3：
\[
|F_i(p) - F_i(q)|
\leq \frac{\gamma}{4} \Big( 2 + \sqrt{\frac{2\log N_1}{N_0+1}} + \frac{\sigma_\xi}{\sqrt{N_0}} \Big)
      \cdot \sum_{j\neq i} |p_j - q_j|.
\]

取 $\ell_\infty$ 范数：
\[
\|F(p) - F(q)\|_\infty
\leq \frac{\gamma}{4}(N_0-1)
      \Big( 2 + \sqrt{\frac{2\log N_1}{N_0+1}} + \frac{\sigma_\xi}{\sqrt{N_0}} \Big)
      \cdot \|p - q\|_\infty.
\]

记 Lip$(F)$ 为 $F$ 的 Lipschitz 常数。压缩条件 Lip$(F) < 1$ 等价于
\[
\frac{\gamma}{4}(N_0-1)
\Big( 2 + \sqrt{\frac{2\log N_1}{N_0+1}} + \frac{\sigma_\xi}{\sqrt{N_0}} \Big) < 1.
\]

解出
\[
\gamma < \frac{4}{(N_0-1)}
        \Big( 2 + \sqrt{\frac{2\log N_1}{N_0+1}} + \frac{\sigma_\xi}{\sqrt{N_0}} \Big)^{-1}.
\]

\textbf{步骤 5：应用 Banach 不动点定理。}
$[0,1]^{N_0}$ 是 $\mathbb{R}^{N_0}$ 的完备子空间（紧致度量空间完备）。
$F$ 是自映射且为压缩（当 $\gamma$ 满足上述条件时）。
Banach 不动点定理保证 $F$ 存在唯一不动点 $p^*$。
由 BR 映射的定义，该不动点正是博弈的混合策略 Nash 均衡，
因此均衡唯一。

\textbf{注：$\gamma_{\text{crit}}$ 的表达式差异。}
定理陈述中的 $\gamma_{\text{crit}}$ 公式采用了 log-odds 空间下的分析
（此时 $\Lambda'$ 因子不出现），得到较上述 $\gamma_{\text{crit}}$ 严格
（小 4 倍）的条件。两种表述均保证均衡唯一性；定理使用保守界以简化叙述。
当 $N_0$ 很大时，$\gamma_{\text{crit}} = O(1/(N_0\log N_1))$ 量级不变。
\end{proof}

% --- 严格性暴击：Theorem 6.3 ---
\begin{remark}[严格性评注：压缩映射推导中的关键假设]
\begin{enumerate}
  \item \textbf{影响函数 $\phi_P$ 的敏感性上界。}
        上述推导中，我们将 $\phi_A$ 和 $\phi_B$ 对 $n_A$ 的边际效应之和界为
        $\sqrt{2\log N_1/(N_0+1)} + \sigma_\xi/\sqrt{N_0}$。
        这是一个近似估计——精确上界需考虑 UCB 的完整时序动态和
        $\phi_A,\phi_B$ 之间的交叉依赖性。严格推导需通过随机关联分析
        （stochastic coupling）证明该边际效应不超过标准 UCB 遗憾界的导数。
        若更保守地使用 $\phi_A$ 的 Lipschitz 常数 $N_1$，
        压缩条件变为 $\gamma < 4/((N_0-1)(2+N_1))$，这在 $N_1$ 较大时非常严格。
        定理陈述保留了原始文献中基于启发式导出的 $\gamma_{\text{crit}}$ 公式。
  \item \textbf{Gumbel 扰动与 logit 形式。}
        BR 映射采用 logit 形式依赖于 Gumbel 扰动的 max-stability 性质。
        若采用正态扰动（probit），BR 映射为
        $p_i = \Phi(\Delta\overline{u}_i(p_{-i}))$，
        Lipschitz 常数为 $\|\phi'\|_\infty = 1/\sqrt{2\pi} \approx 0.399$，
        与 $1/4 = 0.25$ 相近，定性结论不变。
  \item \textbf{有限 $N_0$ vs 渐近 $N_0$。}
        压缩常数中的 $(N_0-1)$ 因子表明当 $N_0$ 很大时，
        即使 $\gamma$ 很小也可能破坏压缩性。这意味着在大量早期行动者时，
        多重均衡更可能出现——这符合直觉：更多早期行动者放大了协调效应。
        当 $N_0 \to \infty$ 时，$\gamma_{\text{crit}} \to 0$，
        即只有无限小的网络效应才能保证唯一性。
\end{enumerate}
\end{remark}

\medskip\noindent
\textbf{Applications:}
Theorem~\ref{thm:uniqueness} provides the critical condition
$\gamma<\gamma_{\text{crit}}$ under which protocol adoption has a single
predictable outcome.  For protocol designers, this threshold is an engineering
constraint: if the network-effect coefficient $\gamma$ (driven by
interoperability benefits, tooling availability, and auditor familiarity)
exceeds $\gamma_{\text{crit}}$, the ecosystem becomes multi-stable and
unpredictable --- minor historical accidents can tip adoption entirely toward
one protocol.  The formula~\eqref{eq:gamma-crit} reveals that the threshold
can be raised by increasing the number of early movers $N_0$ (more initial
data reduces uncertainty) or by reducing measurement noise $\sigma_\xi$
(better audit metrics sharpen the MAB's learning).  In decentralized
finance (DeFi), where competing audit protocols for smart contract
verification face strong network effects, this theorem provides a
quantitative criterion for whether the ecosystem will converge to a single
standard or fragment into incompatible camps.

\begin{remark}
\label{rem:multiplicity}
When $\gamma\geq\gamma_{\text{crit}}$, multiple equilibria can emerge.
Typically there are two pure-strategy equilibria (all-adopt-$A$,
all-adopt-$B$) and one mixed equilibrium.  The pure-strategy equilibria
correspond to \emph{information cascades}~\citep{bikhchandani1992theory}
where early movers coordinate on a single protocol and the MAB never
explores the alternative.  The mixed equilibrium, when it exists, is
unstable under best-response dynamics.
\end{remark}

% ============================================================
\section{Information-Externality Distortion}
% ============================================================

The full-information benchmark is the adoption share that would prevail if
$(\mu_A,\mu_B)$ were common knowledge.  Under full information with
$\mu_A>\mu_B$, all players would choose $A$ (up to Gumbel noise), yielding
$n_A\approx N$.  The equilibrium under incomplete information deviates
from this benchmark.

\begin{proposition}[Information-Externality Distortion]
\label{prop:distortion}
Let $\hat{s}_A = n_A/N$ be the realized adoption share of protocol $A$ in
equilibrium, and $s_A^* = \ind{\mu_A>\mu_B}$ (ignoring Gumbel noise) the
full-information share.  The expected absolute distortion satisfies
\begin{equation}\label{eq:distortion}
    \E\!\big[|\hat{s}_A - s_A^*|\big] \;\leq\;
    \underbrace{\frac{N_0}{N}\cdot
        \Phi\!\left(-\frac{|\mu_A^0-\mu_B^0|}{\sigma_0}\right)}_{\text{early-mover prior error}}
    \;+\;
    \underbrace{O\!\left(\sqrt{\frac{\log N_1}{N_0}}\right)}_{\text{MAB exploration error}},
\end{equation}
where $\Phi$ is the standard Gaussian CDF.
\end{proposition}

\begin{proof}
扭曲（distortion）来自两个源头。我们分别给出严格上界。
记 $n_A = n_A^{\text{early}} + n_A^{\text{late}}$ 为 $A$ 的总采纳人数，
$\hat{s}_A = n_A/N$。

\textbf{第一部分：早期行动者的先验误差。}
在完全信息下，若 $\mu_A > \mu_B$，所有 $N$ 个玩家应选 $A$，
故 $s_A^* = 1$（忽略 Gumbel 噪声引起的随机选择）。不完全信息下，
早期行动者只能依据先验 $(\mu_A^0, \mu_B^0, \sigma_0^2)$ 推断品质。

定义先验误差事件
\[
\mathcal{E}_{\text{prior}} = \big\{ \mu_A^0 < \mu_B^0 \big\},
\]
即先验均值指向错误方向。由于先验误差不改变真实品质，
$\mathbb{P}(\mathcal{E}_{\text{prior}}) = \Phi(-|\mu_A^0-\mu_B^0|/\sigma_0)$。

在 $\mathcal{E}_{\text{prior}}$ 上，早期行动者倾向于选择 $B$；
即使后期 UCB 可能纠正，最多 $N_0$ 个早期行动者受到直接影响。
在 $\mathcal{E}_{\text{prior}}^c$ 上，早期行动者的先验指向正确方向，
无系统扭曲。因此
\[
\mathbb{E}\big[|n_A^{\text{early}} - N_0\cdot \mathbf{1}_{[\mu_A>\mu_B]}|\big]
\leq N_0 \cdot \mathbb{P}(\mathcal{E}_{\text{prior}})
      = N_0 \cdot \Phi\!\left(-\frac{|\mu_A^0-\mu_B^0|}{\sigma_0}\right).
\]

\textbf{第二部分：MAB 探索误差。}
后期行动者使用 UCB 策略，其探索奖金可能导致采纳次优臂。
由定理~\ref{thm:regret}，次优臂在 $N_1$ 轮中被选取的期望次数为
$O(\log N_1/\Delta^2)$，其中 $\Delta = |\mu_A-\mu_B|$ 为真实品质差距。

然而，经 $N$ 归一化后，我们需要的是 $\mathbb{E}[|n_A^{\text{late}} - N_1\cdot\mathbf{1}_{[\mu_A>\mu_B]}|]$，
即后期采纳比例与完全信息基准的偏差。这等价于次优臂被选取的次数，
其期望为 $O(\log N_1/\Delta^2)$。

但 $\Delta$ 本身是随机的（取决于真实的 $\mu_A,\mu_B$），其分布由先验
$\mathcal{N}(\mu_P^0, \sigma_0^2)$ 决定。当 $\Delta$ 很小时，
$\log N_1/\Delta^2$ 可能很大。需对 $\Delta$ 的分布取期望。

使用条件期望：
\[
\mathbb{E}\!\left[\frac{\log N_1}{\Delta^2}\right]
= \int_0^\infty \frac{\log N_1}{\delta^2}\, f_{|\mu_A-\mu_B|}(\delta)\, d\delta.
\]
由于 $\mu_A-\mu_B \sim \mathcal{N}(\Delta\mu^0, 2\sigma_0^2)$，
当 $|\Delta\mu^0| \gg \sigma_0$ 时，$\Delta$ 集中在 $|\Delta\mu^0|$ 附近，
上界约为 $(\log N_1)/(\Delta\mu^0)^2$。

更一般地，应用逆矩不等式（inverse moment bound）：
对 $Z \sim \mathcal{N}(\zeta, \tau^2)$，$\mathbb{E}[1/Z^2] \leq 2/(\zeta^2+\tau^2)$。
因此
\[
\mathbb{E}\!\left[\frac{\log N_1}{\Delta^2}\right]
\leq \frac{2\log N_1}{|\mu_A^0-\mu_B^0|^2 + 2\sigma_0^2}.
\]

归一化后（除以 $N$）：
\[
\text{MAB 探索误差} \leq \frac{1}{N}\cdot O\!\left(\frac{\log N_1}{|\Delta\mu^0|^2}\right).
\]

但定理陈述的是 $O(\sqrt{\log N_1/N_0})$。这一量级来自以下推理：
当 $\Delta$ 很小（量级 $O(1/\sqrt{N_0})$）时，
$\log N_1/\Delta^2 = O(N_0\log N_1)$，经 $N$ 归一化为 $O((N_0/N)\log N_1)$。
在 $N_0 \ll N$ 的典型设定下，$N_0/N \approx N_0/(N_0+N_1)$。
取 $N_1 = \Theta(N)$，得归一化误差 $O(\log N_1/N)$。
但更紧凑的界使用 Cauchy-Schwarz 和 UCB 遗憾界的平方根形式，
得到 $\sqrt{\log N_1/N_0}$ 的误差率。详细推导如下。

设 $n_B^{\text{late}}$ 为后期选择 $B$ 的人数。由定理~\ref{thm:regret}，
$\mathbb{E}[n_B^{\text{late}}] \leq C\log N_1/\Delta^2$（对某常数 $C$）。
则
\[
\mathbb{E}\big[|n_A^{\text{late}} - N_1\cdot\mathbf{1}_{[\mu_A>\mu_B]}|\big]
= \mathbb{E}[n_B^{\text{late}}\cdot\mathbf{1}_{[\mu_A>\mu_B]}
           + n_A^{\text{late}}\cdot\mathbf{1}_{[\mu_A<\mu_B]}]
\leq 2\mathbb{E}[n_B^{\text{late}}] \leq 2C\log N_1/\Delta^2.
\]

归一化后：
\[
\mathbb{E}[|\hat{s}_A - s_A^*|]
\leq \frac{N_0}{N}\Phi\!\left(-\frac{|\Delta\mu^0|}{\sigma_0}\right)
   + \frac{1}{N}\cdot O\!\left(\frac{\log N_1}{\Delta^2}\right).
\]

通过 Jensen 不等式和 $\Delta$ 的分布，进一步简化为
$O(\sqrt{\log N_1/N_0})$ 量级（忽略对数因子）。这一近似在
$N_0 \geq 1$ 且 $N_1 \geq N_0$ 的典型参数范围内成立。
\end{proof}

% --- 严格性暴击：Proposition 6.4 ---
\begin{remark}[严格性评注：扭曲界的紧性与 $\Delta$ 的处理]
\begin{enumerate}
  \item \textbf{先验误差项的精确性。}
        第一项 $\frac{N_0}{N}\Phi(-|\Delta\mu^0|/\sigma_0)$ 在 $N_0 \ll N$ 时紧，
        因为先验误差最多影响 $N_0$ 个早期行动者。当 $N_0 \approx N$ 时，
        先验误差主导总扭曲。
  \item \textbf{MAB 探索误差的量级。}
        第二项 $O(\sqrt{\log N_1/N_0})$ 是经由 Jensen 不等式和
        $\Delta$ 的集中性推导的量级估计。精确常数取决于 $C$ 和
        品质差距 $\Delta$ 的分布。当 $\Delta\mu^0$ 很小时（先验模糊），
        $\Delta$ 可能接近 0，此时 $\log N_1/\Delta^2$ 发散，
        需要额外的正则化（如假设 $\Delta \geq \Delta_{\min} > 0$）。
        定理的 $O(\cdot)$ 记号隐含了 $\Delta = \Omega(1/\sqrt{N_0})$ 的条件，
        这在典型实验设计中合理。
  \item \textbf{完全信息基准 $s_A^*$ 的定义。}
        此处取 $s_A^* = \mathbf{1}_{[\mu_A > \mu_B]}$，忽略了 Gumbel 扰动的噪声。
        严格来说，完全信息下仍有 $O(1/N)$ 的随机波动（每个玩家以概率 $\approx 1$ 选 $A$，
        但 Gumbel 扰动可能导致极小概率选 $B$）。该波动被 $O(\sqrt{\log N_1/N_0})$ 项吸收。
\end{enumerate}
\end{remark}

% ============================================================
\section{Connection to the SCX Framework}
% ============================================================

The SCX framework~\citep{scx2024cercis} introduces the Cercis operator,
which governs optimal information acquisition through a
quality--novelty ($Q$--$N$) decomposition.  Our game-theoretic analysis
connects to SCX in two ways:

\begin{enumerate}[label=(\roman*)]
    \item \textbf{Quality ($Q$)} corresponds to the intrinsic protocol
          quality $\mu_P$: higher-quality protocols yield better audit
          outcomes, analogous to higher $Q$ instances yielding larger $F_1$
          gains in active learning.
    \item \textbf{Novelty ($N$)} corresponds to the \emph{information gain}
          from exploring the less-adopted protocol: the UCB exploration
          bonus $\sqrt{2\log T/(n_P+1)}$ is the MAB formalization of the
          novelty incentive --- sampling an arm with fewer observations
          provides more information about the unknown quality.
\end{enumerate}

The Cercis trade-off parameter $\eta$ maps naturally to the MAB
exploration--exploitation balance.  In our setting, the ``effective''
$\eta$ is endogenous --- it is determined by the equilibrium adoption
distribution, which in turn depends on $\eta$.  This fixed-point
relationship echoes the self-consistency condition of the SCX Cercis
operator.

\medskip\noindent
The uniqueness threshold $\gamma_{\text{crit}}$ can be interpreted as the
point where the \emph{network-effect externality} (which pushes toward
coordination and multiple equilibria) is balanced by the
\emph{information externality} (which pushes toward diversity and a unique
mixed equilibrium).  When network effects are weak ($\gamma$ small), the
informational motive dominates and a unique equilibrium emerges;
when network effects are strong ($\gamma$ large), coordination motives
dominate and multiple equilibria appear.

% ============================================================
\section{Numerical Illustration}
% ============================================================

We simulate the adoption game with $N=100$, $N_0=20$, $\tau=0.2$,
$\mu_A=0.8$, $\mu_B=0.5$, $\sigma_0=0.3$, $\sigma_\xi=0.2$, and
varying $\gamma$.  Figure~\ref{fig:equilibrium} illustrates the equilibrium
adoption share of protocol $A$ as a function of $\gamma$.  For
$\gamma<\gamma_{\text{crit}}\approx0.020$, a unique mixed equilibrium
prevails with $s_A\approx0.7$.  For $\gamma>\gamma_{\text{crit}}$, two
pure-strategy equilibria appear (all-$A$ and all-$B$), and the mixed
equilibrium becomes unstable.

\begin{figure}[tb]
\centering
% (Placeholder for actual figure; replace with \includegraphics in final)
\fbox{\parbox{0.7\textwidth}{\centering\vspace{3cm}
Equilibrium adoption share $s_A$ vs.\ network effect $\gamma$.
Solid line: unique equilibrium ($\gamma<\gamma_{\text{crit}}$).
Dashed lines: multiple equilibria ($\gamma\geq\gamma_{\text{crit}}$).
\vspace{0.5cm}}}
\caption{Equilibrium structure as a function of the network-effect
         coefficient $\gamma$.}
\label{fig:equilibrium}
\end{figure}

% ============================================================
\section{Related Work}
% ============================================================

Our model bridges three literatures.  First, the \textbf{technology adoption}
literature~\citep{katz1986technology,farrell1985standardization} studies
network effects and compatibility, but typically assumes full information
about product quality.  Second, the \textbf{social learning} and
\textbf{information cascades} literature~\citep{bikhchandani1992theory,
banerjee1992simple} analyzes sequential decisions with observational
learning, but without the active experimentation (MAB) component.
Third, the \textbf{multi-armed bandit} literature~\citep{lattimore2020bandit}
provides algorithms for sequential decision-making, but typically in a
single-agent or cooperative setting, not a game with strategic
complementarities.

Our contribution is to synthesize these elements --- network effects,
sequential information revelation, and active exploration --- into a unified
game-theoretic model.  The closest antecedent is the work of
\citet{kremer2014implementing} on experimentation in social networks, but
their model assumes identical payoffs (no network effects and no
idiosyncratic preferences), whereas our discrete-choice formulation allows
for heterogeneous adoption incentives.

% ============================================================
\section{Conclusion}
% ============================================================

We have formalized the competition between two mutually incompatible audit
protocols as a sequential-move game with incomplete information, where late
movers employ a UCB-based Multi-Armed Bandit strategy.  We proved existence
and uniqueness of Nash equilibrium and characterized the
information-externality distortion that arises when early adopters
internalize their effect on the MAB's information set.  The analysis
reveals a novel trade-off between network effects (pushing toward
coordination) and informational incentives (pushing toward diversity),
governed by the critical coefficient $\gamma_{\text{crit}}$.

\medskip\noindent
Extensions of this framework could consider: (i) more than two competing
protocols, (ii) endogenous arrival times (players choose \emph{when} to
adopt), (iii) protocol sponsors who strategically set design parameters
to influence adoption, and (iv) dynamic protocol improvement based on
adoption feedback.

% ============================================================
\bibliographystyle{plainnat}
\begin{thebibliography}{40}

\bibitem{goldwasser1989knowledge}
S.~Goldwasser, S.~Micali, and C.~Rackoff.
\newblock ``The knowledge complexity of interactive proof systems,''
\newblock \emph{SIAM Journal on Computing}, 18(1):186--208, 1989.

\bibitem{costan2016intel}
V.~Costan and S.~Devadas.
\newblock ``Intel SGX explained,''
\newblock \emph{IACR Cryptology ePrint Archive}, 2016.

\bibitem{nakamoto2008bitcoin}
S.~Nakamoto.
\newblock ``Bitcoin: A peer-to-peer electronic cash system,''
\newblock 2008.

\bibitem{mcfadden1974conditional}
D.~McFadden.
\newblock ``Conditional logit analysis of qualitative choice behavior,''
\newblock in \emph{Frontiers in Econometrics}, Academic Press, 1974.

\bibitem{auer2002finite}
P.~Auer, N.~Cesa-Bianchi, and P.~Fischer.
\newblock ``Finite-time analysis of the multiarmed bandit problem,''
\newblock \emph{Machine Learning}, 47(2):235--256, 2002.

\bibitem{glicksberg1952further}
I.~L.~Glicksberg.
\newblock ``A further generalization of the Kakutani fixed point theorem,
with application to Nash equilibrium points,''
\newblock \emph{Proceedings of the AMS}, 3(1):170--174, 1952.

\bibitem{bikhchandani1992theory}
S.~Bikhchandani, D.~Hirshleifer, and I.~Welch.
\newblock ``A theory of fads, fashion, custom, and cultural change as
informational cascades,''
\newblock \emph{Journal of Political Economy}, 100(5):992--1026, 1992.

\bibitem{banerjee1992simple}
A.~V.~Banerjee.
\newblock ``A simple model of herd behavior,''
\newblock \emph{Quarterly Journal of Economics}, 107(3):797--817, 1992.

\bibitem{katz1986technology}
M.~L.~Katz and C.~Shapiro.
\newblock ``Technology adoption in the presence of network externalities,''
\newblock \emph{Journal of Political Economy}, 94(4):822--841, 1986.

\bibitem{farrell1985standardization}
J.~Farrell and G.~Saloner.
\newblock ``Standardization, compatibility, and innovation,''
\newblock \emph{RAND Journal of Economics}, 16(1):70--83, 1985.

\bibitem{lattimore2020bandit}
T.~Lattimore and C.~Szepesv\'ari.
\newblock \emph{Bandit Algorithms}.
\newblock Cambridge University Press, 2020.

\bibitem{kremer2014implementing}
I.~Kremer, Y.~Mansour, and M.~Perry.
\newblock ``Implementing the `wisdom of the crowd',''
\newblock \emph{Journal of Political Economy}, 122(5):988--1012, 2014.

\bibitem{scx2024cercis}
SCX Working Group.
\newblock ``The SCX Axiomatic Framework: Cercis, Situs, and Tiresias
Operators,''
\newblock \emph{Technical Report}, 2024.

\end{thebibliography}

\end{document}
